% ============================================================
% CAPITOLO 16 — Pattern Avanzati e Oltre
% ============================================================
\chapter{Pattern Avanzati e Oltre}\index{microservizi}

\section*{Cosa Imparerai}
Questo capitolo non ha un progetto da costruire. È una mappa per il futuro:
\begin{itemize}
  \item Come applicare il metodo 0-code a progetti enterprise
  \item Pattern Multi-Agente avanzati (orchestrazione reale)
  \item Lavorare con codice legacy
  \item Architetture a microservizi
  \item I limiti onesti del paradigma 0-code
  \item Come continuare a crescere
\end{itemize}

\textbf{Tempo di lettura}: 30--40 minuti

% ---------------------------------------------------------
\section{Scalare il Metodo}

\subsection{Da progetto singolo a team}

Tutto quello che hai fatto finora era per un singolo sviluppatore. In un team, il metodo 0-code richiede adattamenti.

\textbf{Il \texttt{\_CONTEXT.md} diventa un contratto di team.}

\begin{lstlisting}[language={},title={Contesto di team}]
# Progetto: E-Commerce Platform

## Team
- Backend Team: API, database, autenticazione
  (contesto: ./backend/_CONTEXT.md)
- Frontend Team: React SPA, mobile app
  (contesto: ./frontend/_CONTEXT.md)
- Ops Team: infrastruttura, CI/CD, monitoring
  (contesto: ./infra/_CONTEXT.md)

## Regole di Collaborazione
- Ogni modifica agli endpoint API deve aggiornare
  il _CONTEXT.md di root
- Le PR che cambiano il contratto API richiedono
  approvazione da entrambi i team
- Il _CONTEXT.md e versionato in git
\end{lstlisting}

\subsection{Convenzione dei contesti gerarchici}

\begin{lstlisting}[language={}]
progetto/
+-- _CONTEXT.md              <- Architettura globale
+-- backend/
|   +-- _CONTEXT.md          <- Stack, modelli backend
|   +-- SKILL.md             <- Competenze (Prisma, Auth)
|   +-- services/
|       +-- payment/
|           +-- _CONTEXT.md  <- Microservizio pagamenti
+-- frontend/
|   +-- _CONTEXT.md          <- Convenzioni React
|   +-- SKILL.md             <- Design system
+-- mobile/
    +-- _CONTEXT.md          <- Convenzioni Flutter
    +-- SKILL.md             <- Pattern nativi
\end{lstlisting}

L'IA legge il contesto \textbf{dal più specifico al più generale}: prima il \texttt{\_CONTEXT.md} della directory corrente, poi quello del padre, poi quello della root. Questo è il \textbf{Progressive Disclosure applicato all'architettura}.

% ---------------------------------------------------------
\section{Pattern Multi-Agente Avanzati}

Nel Capitolo~10 hai usato il pattern Planner-Generator-Evaluator cambiando manualmente il ruolo dell'IA. Nei sistemi avanzati, puoi automatizzare questo flusso.

\subsection{Orchestrazione con file di istruzioni}

Crea file \texttt{.instructions.md} o \texttt{.agent.md} che definiscono comportamenti specifici:

\begin{lstlisting}[language={},title={.github/copilot-instructions.md}]
# Istruzioni per lo sviluppo di questo progetto

Quando ricevi una richiesta di nuova funzionalita:
1. PRIMA analizza l'impatto leggendo tutti i _CONTEXT.md
2. POI proponi un piano (file da modificare, ordine, rischi)
3. ASPETTA conferma prima di implementare
4. Dopo l'implementazione, self-review su sicurezza
5. Aggiorna i _CONTEXT.md coinvolti
6. Aggiorna PROGRESS.md
\end{lstlisting}

\subsection{Multi-Tool con MCP}

Nel Capitolo~7 hai usato MCP per PostgreSQL. In scenari avanzati, puoi connettere molteplici tool:

\begin{lstlisting}[language={},title={Configurazione MCP multi-tool}]
{
  "mcpServers": {
    "postgres": {
      "command": "npx",
      "args": ["-y",
        "@modelcontextprotocol/server-postgres",
        "postgresql://..."]
    },
    "filesystem": {
      "command": "npx",
      "args": ["-y",
        "@modelcontextprotocol/server-filesystem",
        "/path/to/docs"]
    },
    "github": {
      "command": "npx",
      "args": ["-y",
        "@modelcontextprotocol/server-github"],
      "env": {
        "GITHUB_PERSONAL_ACCESS_TOKEN": "${GITHUB_TOKEN}"
      }
    }
  }
}
\end{lstlisting}

Con questa configurazione, l'IA può leggere lo schema del database, consultare documentazione locale e creare issue e PR su GitHub --- tutto senza uscire dall'editor.

\subsection{L'Ecosistema MCP nel 2026}

Nel 2026 lo standard MCP\index{MCP!ecosistema 2026} ha superato i 97~milioni di installazioni, guadagnandosi l'appellativo di ``USB-C per l'Intelligenza Artificiale''. Oggi esiste un server MCP per quasi ogni servizio. Ecco le categorie fondamentali per lo sviluppo full-stack:

\begin{center}
\begin{small}
\begin{tabularx}{\textwidth}{l l X}
\toprule
\textbf{Categoria} & \textbf{Server MCP} & \textbf{A cosa serve} \\
\midrule
Documentazione & Context7 & Accesso versionato alla documentazione ufficiale delle librerie. L'agente consulta le API esatte anziché basarsi sui dati di addestramento. \\
Sandbox codice & E2B & Ambiente cloud isolato per eseguire script senza rischi sulla macchina locale. \\
Infrastruttura dati & Supabase / Neon & Gestione database di produzione con policy di sicurezza a livello di riga. \\
Deploy e CI/CD & Vercel / Docker Hub & L'agente legge log di build, analizza errori di deploy, gestisce container. \\
Automazione & Zapier / Composio & Trigger e azioni per servizi aziendali (Slack, Jira, CRM). \\
Design & Figma & Esportazione delle specifiche di design direttamente nel contesto dell'agente. \\
\bottomrule
\end{tabularx}
\end{small}
\end{center}

\begin{warnbox}
Ogni server MCP attivo consuma una porzione della finestra di contesto dell'agente. Non installarne~20 ``per sicurezza'': seleziona solo quelli necessari al progetto corrente. Una buona regola: \textbf{3--5 server MCP} per progetto, scelti in base allo stack tecnologico.
\end{warnbox}

% ---------------------------------------------------------
\section{Lavorare con Codice Legacy}\index{codice legacy}\index{Reverse Context Engineering}

Il 90\% del lavoro di sviluppo reale è su codice esistente, non su progetti nuovi. Il metodo 0-code funziona anche qui.

\subsection{Strategia: Reverse Context Engineering}

\begin{lstlisting}[language={}]
Ho ereditato un progetto PHP/Laravel senza documentazione.
Analizza la struttura e genera un _CONTEXT.md che descriva:
1. Lo scopo dell'applicazione (dedotto da modelli e route)
2. Lo stack tecnologico (versioni, dipendenze)
3. L'architettura (pattern usati, struttura directory)
4. I modelli dati e le loro relazioni
5. Gli endpoint API (se presenti)
6. Le variabili d'ambiente usate
7. Come avviare il progetto in locale
\end{lstlisting}

L'IA analizza il codice e genera il contesto che non esisteva. Da quel momento in poi, puoi lavorare in 0-code su quel progetto.

\subsection{Refactoring incrementale}

\begin{lstlisting}[language={}]
Il progetto usa jQuery per il frontend.
Voglio migrare gradualmente a React, un componente alla volta.

Piano:
1. Identifica i componenti jQuery piu isolati
2. Per ognuno, crea l'equivalente React
3. Approccio "strangler fig": React e jQuery coesistono
4. Procedi un componente alla volta

Inizia dall'analisi: quali componenti jQuery sono
i migliori candidati per la prima migrazione?
\end{lstlisting}

% ---------------------------------------------------------
\section{Microservizi}\index{microservizi}

\subsection{Quando passare ai microservizi}

\textbf{Non} passare ai microservizi perché ``è moderno''. Passa solo quando:
\begin{itemize}
  \item Il monolite ha tempi di deploy troppo lunghi ($>$ 30 minuti)
  \item Parti diverse hanno requisiti di scaling diversi
  \item Team indipendenti devono poter fare deploy senza coordinarsi
\end{itemize}

\subsection{Contesto per microservizi}

Ogni microservizio ha il suo \texttt{\_CONTEXT.md}:

\begin{lstlisting}[language={}]
platform/
+-- _CONTEXT.md              <- Architettura globale
+-- gateway/
|   +-- _CONTEXT.md          <- API Gateway, routing
+-- auth-service/
|   +-- _CONTEXT.md          <- JWT, OAuth, utenti
+-- notes-service/
|   +-- _CONTEXT.md          <- CRUD note, categorie
+-- notification-service/
|   +-- _CONTEXT.md          <- Email, push notification
+-- shared/
    +-- _CONTEXT.md          <- Librerie condivise, DTOs
    +-- SKILL.md             <- Convenzioni condivise
\end{lstlisting}

Il \texttt{\_CONTEXT.md} della root definisce i \textbf{contratti tra servizi}:

\begin{center}
\begin{tabularx}{\textwidth}{l l l X}
\toprule
\textbf{Da} & \textbf{A} & \textbf{Metodo} & \textbf{Descrizione} \\
\midrule
Gateway & Auth & HTTP & Verifica token \\
Gateway & Notes & HTTP & Proxy richieste note \\
Notes & Notification & Event (Redis) & Nuova nota $\rightarrow$ notifica \\
\bottomrule
\end{tabularx}
\end{center}

\begin{center}
\begin{tabularx}{\textwidth}{l l l X}
\toprule
\textbf{Evento} & \textbf{Producer} & \textbf{Consumer} & \textbf{Payload} \\
\midrule
note.created & Notes & Notification & \texttt{\{ userId, noteId, title \}} \\
user.deleted & Auth & Notes, Notif. & \texttt{\{ userId \}} \\
\bottomrule
\end{tabularx}
\end{center}

% ---------------------------------------------------------
\section{I Limiti Onesti del Paradigma 0-Code}

\subsection{Dove funziona bene}

\begin{center}
\begin{tabularx}{\textwidth}{l X}
\toprule
\textbf{Scenario} & \textbf{Perché} \\
\midrule
CRUD applications & Pattern ben noti, l'IA eccelle \\
REST API & Standard consolidato, tanta documentazione \\
Frontend React/Flutter & Componenti modulari, pattern ripetitivi \\
Integrazione servizi & OAuth, pagamenti, email --- pattern standard \\
Testing & I test sono pattern ripetitivi per natura \\
DevOps/CI-CD & Configurazioni dichiarative, template noti \\
\bottomrule
\end{tabularx}
\end{center}

\subsection{Dove fatica}

\begin{center}
\begin{small}
\begin{tabularx}{\textwidth}{l X X}
\toprule
\textbf{Scenario} & \textbf{Perché} & \textbf{Cosa fare} \\
\midrule
Algoritmi complessi & Codice corretto ma non ottimizzato & Specifica la complessità richiesta \\
Real-time/WebSocket & Meno esempi, più errori & Fornisci più contesto \\
Sistemi distribuiti & Troppi edge case & Pianifica tu, delega l'implementazione \\
Performance critiche & Ottimizza per correttezza & Profila, chiedi di ottimizzare il bottleneck \\
Domini specializzati & Finanza, bioinformatica & SKILL.md con conoscenza di dominio \\
Creatività pura & Game design, UX innovativa & L'IA replica pattern, non li inventa \\
\bottomrule
\end{tabularx}
\end{small}
\end{center}

\subsection{Regola pratica}

\begin{notebox}
Se riesci a descrivere cosa vuoi in 2--3 paragrafi di testo chiaro, l'IA lo può implementare. Se non riesci a descriverlo chiaramente, probabilmente non hai ancora capito abbastanza il problema --- e in quel caso nemmeno tu lo implementeresti bene.
\end{notebox}

% ---------------------------------------------------------
\section{Il Futuro del Paradigma}

\subsection{Tendenze in corso}

\begin{enumerate}
  \item \textbf{Modelli sempre più capaci}: Claude, GPT e altri migliorano velocemente. Quello che oggi richiede 5~iterazioni, domani ne richiederà~2.
  \item \textbf{Contesto più grande}: Le finestre di contesto crescono. Progetti interi possono essere analizzati in una singola richiesta.
  \item \textbf{Tool sempre più integrati}: MCP e protocolli simili collegano l'IA direttamente ai sistemi (database, cloud, monitoring).
  \item \textbf{Agenti autonomi}: L'IA non solo genera codice, ma esegue test, deploya e corregge errori. Il tuo ruolo diventa sempre più strategico.
\end{enumerate}

\subsection{Cosa resta costante}

Indipendentemente dall'evoluzione dei modelli, le competenze che hai costruito in questo libro restano valide:
\begin{itemize}
  \item \textbf{Saper descrivere} cosa vuoi (Context Engineering)
  \item \textbf{Saper verificare} che sia corretto (Confidence Tagging)
  \item \textbf{Saper strutturare} il lavoro (ADLC)
  \item \textbf{Saper valutare il rischio} (OWASP, code review)
\end{itemize}

Queste non sono competenze sull'IA. Sono competenze di \emph{ingegneria del software}.

% ---------------------------------------------------------
\section{Vibe Coding e il Dibattito FAAFO}

Nel 2026 il paradigma dominante nella comunità degli sviluppatori ha preso il nome di \textbf{Vibe Coding}\index{Vibe Coding}, un termine coniato da Andrej Karpathy\index{Karpathy, Andrej} e successivamente strutturato in un framework completo da Gene Kim e Steve Yegge~\cite{kim-yegge:vibe-coding}. Il Vibe Coding descrive la pratica di generare software descrivendo l'\textit{atmosfera} o l'intento ad alto livello in linguaggio naturale, affidandosi all'inferenza del modello piuttosto che a specifiche algoritmiche esatte.

Kim e Yegge hanno introdotto il framework \textbf{FAAFO}\index{FAAFO} per governare questa transizione~\cite{kim-yegge:vibe-coding,yegge:tips}:

\begin{center}
\begin{tabularx}{\textwidth}{l l X}
\toprule
\textbf{Lettera} & \textbf{Principio} & \textbf{Significato} \\
\midrule
F & Fast & I tempi di ciclo crollano: prototipi in minuti, non mesi \\
A & Ambitious & Un singolo sviluppatore affronta progetti da team intero \\
A & Autonomous & Lo sviluppatore passa da ``cuoco di linea'' a ``capocuoco'' \\
F & Fun & L'attrito del debugging sintattico scompare \\
O & Optionality & Genera decine di varianti per valutare empiricamente la migliore \\
\bottomrule
\end{tabularx}
\end{center}

\subsection{Il concetto di ``software a perdere''}

L'aspetto più controverso del Vibe Coding è l'idea che il software moderno sia intrinsecamente \textit{scartabile}: quando i requisiti cambiano, anziché mantenere il codice esistente può essere più efficiente istruire l'agente a riscrivere intere sezioni da zero.

\begin{notebox}
Questo principio ha validità per la \textbf{prototipazione rapida} e la validazione di idee. Per i sistemi in produzione con utenti reali, dati persistenti e requisiti di conformità, il metodo strutturato dell'ADLC resta necessario. Il Vibe Coding e l'ADLC non sono in opposizione: sono strumenti per fasi diverse del ciclo di vita del prodotto.
\end{notebox}

\subsection{Come bilanciare i due approcci}

\begin{center}
\begin{small}
\begin{tabularx}{\textwidth}{l X X}
\toprule
\textbf{Fase} & \textbf{Approccio Vibe Coding} & \textbf{Approccio ADLC Strutturato} \\
\midrule
Esplorazione idea & Prompt vago, genera 5 varianti, scegli & Non applicabile --- troppo overhead \\
MVP/Prototipo & Genera tutto, valida con utenti & \texttt{\_CONTEXT.md} minimo, ADLC leggero \\
Produzione & Non sufficiente da solo & ADLC completo, Confidence Tagging, OWASP \\
Manutenzione & Riscrivi le parti obsolete & Context Engineering + versioning \\
\bottomrule
\end{tabularx}
\end{small}
\end{center}

\begin{tipbox}
Il percorso di questo libro ti ha insegnato l'ADLC strutturato perché è il metodo più sicuro e didatticamente solido. Ora che lo padroneggi, puoi permetterti di ``vibrare'' nelle fasi esplorative, sapendo esattamente \textit{quando} tornare al rigore.
\end{tipbox}

% ---------------------------------------------------------
\section{Piattaforme App Builder AI-Native}

Non sempre è necessario orchestrare manualmente backend, frontend e database. Nel 2026 sono emerse piattaforme che astraggono l'intero stack, generando applicazioni complete da una descrizione in linguaggio naturale:

\begin{center}
\begin{tabularx}{\textwidth}{l X}
\toprule
\textbf{Piattaforma} & \textbf{Caratteristica distintiva} \\
\midrule
Base44 & App funzionante in minuti, ideale per tool interni e dashboard \\
Lovable & Prototipazione fulminea di frontend e interfacce \\
Bolt.new & Generazione di web app complete con deploy istantaneo \\
Vybe & Integra agenti IA nell'app finale per task amministrativi continuativi \\
\bottomrule
\end{tabularx}
\end{center}

\subsection{Quando usare un App Builder vs.\ orchestrazione manuale}

\begin{center}
\begin{small}
\begin{tabularx}{\textwidth}{X X}
\toprule
\textbf{App Builder AI-Nativo} & \textbf{Orchestrazione Manuale (questo libro)} \\
\midrule
Priorità assoluta: velocità di consegna & Controllo granulare sull'infrastruttura \\
MVP, validazione di mercato & Personalizzazioni estreme, sistemi regolamentati \\
Tool interni con logica business standard & Integrazione con sistemi legacy \\
Budget limitato, team di 1 persona & Scalabilità e performance critiche \\
\bottomrule
\end{tabularx}
\end{small}
\end{center}

\begin{notebox}
Comprendere \textit{quando non scrivere codice} --- nemmeno tramite un agente --- è la forma più alta di efficienza del paradigma 0-Code. Ma quando il progetto richiede sicurezza, personalizzazione e controllo, il metodo che hai appreso in 15~capitoli resta la scelta professionale.
\end{notebox}

% ---------------------------------------------------------
\section{Continuare a Crescere}

\subsection{Progetti suggeriti per consolidare}

\begin{center}
\begin{tabularx}{\textwidth}{l X}
\toprule
\textbf{Progetto} & \textbf{Nuova competenza} \\
\midrule
Blog con CMS & Server-side rendering (Next.js), SEO, Markdown parsing \\
Chat real-time & WebSocket, Socket.IO, messaggi in tempo reale \\
E-commerce & Pagamenti (Stripe), carrello, ordini, email transazionali \\
Dashboard analytics & Grafici (Chart.js/D3), aggregazioni SQL, export \\
App IoT & MQTT, dati in streaming, time-series database \\
\bottomrule
\end{tabularx}
\end{center}

\subsection{Per ogni progetto}

\begin{enumerate}
  \item Scrivi il \texttt{\_CONTEXT.md} prima di tutto
  \item Crea il \texttt{SKILL.md} se il dominio è nuovo
  \item Segui le 7~fasi ADLC
  \item Esegui il Confidence Tagging sulle parti critiche
  \item Aggiorna \texttt{PROGRESS.md} tra le sessioni
\end{enumerate}

% ---------------------------------------------------------
\section*{Conclusione}

Quando hai iniziato questo libro, l'idea di costruire un'applicazione completa --- backend, frontend, mobile, autenticazione, database, deploy --- era probabilmente intimidatoria.

Ora l'hai fatto. E il codice sorgente di queste applicazioni non l'hai scritto tu. L'hai \textbf{descritto}.

Lo sviluppo 0-code non è magia. È un metodo disciplinato: contesto preciso, verifica rigorosa, iterazione continua. Gli strumenti cambieranno. I modelli miglioreranno. Ma il metodo --- descrivere, generare, verificare --- è qui per restare.

Buon sviluppo.
